% ========================================
% 附录
% ========================================

\chapter*{附\quad 录}
\addcontentsline{toc}{chapter}{附录}

\section*{附录A 开发环境配置指南}

\subsection*{A.1 Java开发环境}

\begin{enumerate}
\item \textbf{JDK安装}
   \begin{itemize}
   \item 下载地址：https://www.oracle.com/java/technologies/downloads/
   \item 推荐版本：JDK 11 或 JDK 17
   \item 配置环境变量：JAVA\_HOME, PATH
   \end{itemize}

\item \textbf{Maven配置}
   \begin{itemize}
   \item 下载地址：https://maven.apache.org/download.cgi
   \item 配置环境变量：MAVEN\_HOME, PATH
   \item 配置镜像源：阿里云或华为云镜像
   \end{itemize}
\end{enumerate}

\subsection*{A.2 前端开发环境}

\begin{enumerate}
\item \textbf{Node.js安装}
   \begin{itemize}
   \item 下载地址：https://nodejs.org/
   \item 推荐版本：Node.js 16+ 或 18+
   \item 包管理器：npm 或 yarn
   \end{itemize}

\item \textbf{Vue CLI安装}
   \begin{lstlisting}[style=bash]
npm install -g @vue/cli
vue --version
   \end{lstlisting}
\end{enumerate}

\section*{附录B 常用命令参考}

\subsection*{B.1 Docker命令}

\begin{lstlisting}[style=bash]
# 构建镜像
docker build -t image-name .

# 运行容器
docker run -d -p 8080:8080 image-name

# 查看运行中的容器
docker ps

# 停止容器
docker stop container-id

# 查看日志
docker logs container-id
\end{lstlisting}

\subsection*{B.2 Kubernetes命令}

\begin{lstlisting}[style=bash]
# 应用配置
kubectl apply -f deployment.yaml

# 查看Pod状态
kubectl get pods

# 查看服务
kubectl get services

# 查看日志
kubectl logs pod-name

# 进入容器
kubectl exec -it pod-name -- /bin/bash
\end{lstlisting}

\section*{附录C 配置文件模板}

\subsection*{C.1 Docker Compose配置}

\begin{lstlisting}[style=yaml]
version: '3.8'

services:
  app:
    build: .
    ports:
      - "8080:8080"
    environment:
      - SPRING_PROFILES_ACTIVE=prod
    depends_on:
      - database
      - redis

  database:
    image: postgres:13
    environment:
      POSTGRES_DB: waterconservancy
      POSTGRES_USER: admin
      POSTGRES_PASSWORD: password
    volumes:
      - db_data:/var/lib/postgresql/data

  redis:
    image: redis:6
    ports:
      - "6379:6379"

volumes:
  db_data:
\end{lstlisting}

\subsection*{C.2 Spring Boot配置}

\begin{lstlisting}[style=yaml]
server:
  port: 8080
  servlet:
    context-path: /api

spring:
  datasource:
    url: jdbc:postgresql://localhost:5432/waterconservancy
    username: admin
    password: password
    driver-class-name: org.postgresql.Driver
  
  jpa:
    hibernate:
      ddl-auto: update
    show-sql: false
    properties:
      hibernate:
        dialect: org.hibernate.dialect.PostgreSQLDialect

  redis:
    host: localhost
    port: 6379
    timeout: 2000ms

logging:
  level:
    com.waterconservancy: DEBUG
  pattern:
    console: "%d{yyyy-MM-dd HH:mm:ss} - %msg%n"
    file: "%d{yyyy-MM-dd HH:mm:ss} [%thread] %-5level %logger{36} - %msg%n"
\end{lstlisting}

\section*{附录D 故障排查指南}

\subsection*{D.1 常见问题及解决方案}

\begin{enumerate}
\item \textbf{端口占用问题}
   \begin{lstlisting}[style=bash]
# 查看端口占用
netstat -an | grep :8080
# 或使用
lsof -i :8080

# 杀死进程
kill -9 PID
   \end{lstlisting}

\item \textbf{内存不足}
   \begin{itemize}
   \item 增加JVM堆内存：-Xmx2g
   \item 优化数据库连接池配置
   \item 使用缓存减少数据库访问
   \end{itemize}

\item \textbf{数据库连接失败}
   \begin{itemize}
   \item 检查数据库服务状态
   \item 验证连接字符串和凭据
   \item 确认网络连接和防火墙设置
   \end{itemize}
\end{enumerate}

\subsection*{D.2 性能优化建议}

\begin{enumerate}
\item \textbf{数据库优化}
   \begin{itemize}
   \item 创建合适的索引
   \item 优化SQL查询语句
   \item 使用连接池
   \item 分库分表策略
   \end{itemize}

\item \textbf{应用优化}
   \begin{itemize}
   \item 使用缓存技术（Redis）
   \item 异步处理长时间任务
   \item 代码优化和重构
   \item 合理使用线程池
   \end{itemize}
\end{enumerate}
